% ══════════════════════════════════════════════════════════════════════════════
%  CAPITOLO — LA DASHBOARD INTERATTIVA
%  Etichette: ch:dash, sec:dash_ui, sec:dash_diverge,
%             sec:dash_overdisp, sec:dash_limits
% ══════════════════════════════════════════════════════════════════════════════

\chapter{Una Dashboard Interattiva per l'Esplorazione dei Dati}
\label{ch:dash}

I capitoli precedenti hanno prodotto una sequenza di risultati dipendenti da
intervalli e costanti fisse: una
distribuzione per età, una curva di retention, una tabella di IRR, una
proiezione al 2031. Ognuno risponde a una domanda posta \emph{a priori}.
Questa appendice descrive uno strumento che serve al caso opposto --- rispondere
a domande impreviste e non considerate: Come si
distribuiscono le donazioni dentro un singolo mese del 2020, che aspetto ha la
retention se si sposta la soglia di ``giovane'' da 25 a 30 anni, quanto pesa la
plasmaferesi nel solo bacino femminile. Sono domande legittime e numerose, ed
eseguire ogni analisi richiederebbe più tempo e spazio di quanto a disposizione.

Lo strumento è un'applicazione web in \textbf{R} con il framework
\textbf{Shiny}, che legge lo stesso dataset dei capitoli precedenti e ne
permette l'interrogazione attraverso otto schede tematiche e quattro filtri
globali. Non è un risultato statistico e non aggiunge evidenza: ha semplicemente
il compito di rendere i dati accessibili e leggibili.
% Ha però prodotto, come sottoprodotto della sua
% verifica, un'osservazione metodologica che ha finito per cambiare la specifica
% del \cref{ch:glm}, ed è la ragione principale per cui il capitolo esiste.

% ───────────────────────────────────────────────────────────────────────
\section{Struttura e criterio di progetto}
\label{sec:dash_ui}

Quattro filtri globali --- anno, genere, tipologia, fascia d'età --- agiscono su
tutte le schede, con una sola eccezione motivata: la sotto-scheda sulla
retention ignora il filtro sull'anno, perché restringere a una singola
annualità troncherebbe artificialmente ogni misura di fidelizzazione, che è per
definizione longitudinale.

Le otto schede coprono l'esplorazione descrittiva (\emph{Overview},
\emph{Univariate}, \emph{Time Series}, \emph{Età}), il livello giornaliero
(\emph{Calendar Heatmap}), l'ispezione delle righe (\emph{Data Table}), la
regressione (\emph{Regression}) e le proiezioni (\emph{Forecast}). Il livello
giornaliero è quello che giustifica meglio lo strumento: renderlo in forma
stampata richiederebbe 204 figure, una per mese della serie, mentre nella
dashboard è una scheda sola in cui si naviga. Il gradiente per giorno della
settimana è immediato --- i venerdì sistematicamente più intensi, le domeniche
quasi vuote, l'Epifania visibile come singola cella fredda in un lunedì
altrimenti attivo.

Il criterio di progetto è che \textbf{l'applicazione non calcola nulla a
runtime}: legge artefatti già prodotti. Le uniche stime calcolate al volo sono
i modelli della scheda \emph{Regression}, per la ragione discussa nella
\cref{sec:dash_diverge}.

La scheda \emph{Forecast} legge gli artefatti dell'\cref{ch:forecast}: la
proiezione annuale del trend smorzato (97\,230 donazioni nel sessennio
2026--2031), ripartita sui mesi secondo il profilo stagionale storico.

% ───────────────────────────────────────────────────────────────────────
\section{Dove la dashboard e la tesi divergono}
\label{sec:dash_diverge}

Un'implementazione indipendente degli stessi calcoli, in un linguaggio diverso
e con un percorso di aggregazione diverso, è un controllo sui numeri della
tesi. Dove i due percorsi convergono la fiducia aumenta; dove divergono, la
divergenza è essa stessa informativa --- e si è rivelata la parte più utile.
Due schede producono numeri che \emph{non} coincidono con i capitoli
corrispondenti, e in entrambi i casi la causa non è un errore di calcolo ma una
diversa definizione dell'oggetto misurato.

\paragraph{Retention: due definizioni di ``donatore giovane''.}
Il \cref{ch:retention} classifica come giovane chi ha una data di prima
donazione registrata a un'età fra 18 e 25 anni; la dashboard usa la minima età
\emph{osservata} nel dataset. Per un donatore la cui carriera precede il 2009 i
due criteri lo classificano diversamente, e la dashboard include inoltre le
coorti 2024--2025, che hanno per costruzione una retention osservabile di al
più un anno o due.

Questa divergenza si è rivelata più significativa di quanto sembrasse: è
esattamente la distinzione fra prima donazione osservata e ingresso reale che
la \cref{sec:ret_def} ha poi dovuto affrontare, e che ha portato a rifare le
curve di sopravvivenza del capitolo. La dashboard non ha risolto il problema
--- adotta anch'essa una definizione approssimata --- ma renderlo visibile è
ciò che ha portato a cercarlo.

\paragraph{Regressione: due unità di analisi.}
La scheda \emph{Regression} stima gli stessi due modelli del \cref{ch:glm} ma
su una tabella costruita in modo diverso: aggrega preventivamente a una riga
per donatore e tipologia, contando le donazioni entro l'ambito selezionato dai
filtri, anziché stimare sul panel donatore-anno con gli anni a zero.

La conseguenza è che i coefficienti non sono direttamente confrontabili, e in un
caso non lo sono affatto. Con il filtro anno su ``All'' la risposta della
dashboard è un conteggio \emph{di carriera} anziché annuale, e in assenza di un
termine di esposizione l'età assorbe gli anni di attività accumulati: l'IRR per
anno d'età sale a $1{,}041$. Restringendo il filtro a una singola annualità
scende a $1{,}007$, dello stesso ordine di grandezza del capitolo.

È lo stesso meccanismo diagnosticato nella \cref{sec:glm_effects}: senza un
controllo per anzianità di carriera, il coefficiente dell'età misura
l'anzianità. La dashboard lo ha mostrato prima, e in forma più grossolana, di
quanto il modello del capitolo lo abbia poi quantificato.

% ───────────────────────────────────────────────────────────────────────
\section{Un risultato collaterale: la sovradispersione dipende dall'unità di analisi}
\label{sec:dash_overdisp}

La verifica più utile prodotta dalla dashboard riguarda la scelta della
distribuzione. Confrontando AIC di Poisson e binomiale negativa al variare del
filtro anno si osserva che su tutti gli anni aggregati la binomiale negativa
domina nettamente, ma \textbf{su una singola annualità} --- cioè quando la
risposta è un conteggio annuo non replicato --- il parametro di dispersione
diverge --- la stima non converge, e $\hat{\theta}$ scappa verso l'infinito ---
e l'AIC della binomiale negativa diventa \emph{peggiore} di quello della
Poisson di esattamente due punti, cioè la penalizzazione per il parametro
aggiuntivo e nient'altro. In altri termini la
binomiale negativa collassa sulla Poisson, e la sovradispersione sparisce.

\begin{table}[ht]
\centering
\caption{Confronto AIC nella scheda \emph{Regression} al variare del filtro
anno, ricalcolato sul dataset del \cref{ch:data}. Covariate: età, genere,
tipologia. $\Delta\text{AIC} = \text{AIC}_{\text{Poisson}} -
\text{AIC}_{\text{NB}}$: positivo favorisce la binomiale negativa.}
\label{tab:dash_overdisp}
\begin{tabular}{@{}lrrrrr@{}}
\toprule
\textbf{Filtro anno} & \textbf{$n$} & \textbf{AIC Poisson} & \textbf{AIC NB} &
\textbf{$\Delta$AIC} & \textbf{$\hat{\theta}$ (NB)} \\
\midrule
2015  & 12\,790 & 36\,404{,}4  & 36\,406{,}4  & $-2{,}0$    & diverge \\
2025  & 10\,607 & 29\,895{,}8  & 29\,897{,}8  & $-2{,}0$    & diverge \\
Tutti & 52\,163 & 493\,871{,}1 & 289\,645{,}2 & $+204\,226$ & $1{,}40$ \\
\bottomrule
\end{tabular}
\end{table}

L'osservazione ha sollevato una domanda seria sul \cref{ch:glm}: se la
sovradispersione su cui poggia la scelta della binomiale negativa sparisce
appena si guarda un anno alla volta, quella scelta va motivata diversamente.

\textbf{La domanda ha trovato risposta}, ed è riportata nella
\cref{sec:glm_overdisp}. La sovradispersione che la dashboard misura su tutti
gli anni viene dall'aggregazione longitudinale --- sommando le donazioni di un
donatore su tutta la carriera, senza un termine di esposizione, si sommano
donatori con due anni e donatori con diciassette, e la varianza cresce per una
ragione puramente contabile. Rimossa questa componente, la Poisson sembra
adeguata.

Ma la dashboard, guardando un anno alla volta, guarda anche \emph{soltanto i
donatori attivi in quell'anno}: gli anni a zero non ci sono. È una popolazione
già selezionata, e la sovradispersione che sparisce sparisce per questo.
Materializzando gli anni a zero il rapporto varianza/media passa da $1{,}23$ a
$2{,}45$ e la binomiale negativa torna a dominare con un margine enorme
(\cref{tab:glm_fit}).

L'appendice si chiude quindi con un esito che vale la pena esplicitare: uno
strumento costruito per esplorare i dati ha portato a correggere l'unità di
analisi del modello principale della tesi. È il tipo di ritorno che giustifica
il costo di un'implementazione indipendente.

% ───────────────────────────────────────────────────────────────────────
\section{Limiti}
\label{sec:dash_limits}

L'applicazione non è deployata in nessun ambiente e legge il dataset dal filesystem. 
Renderla accessibile al DIMT richiederebbe un server e soprattutto una valutazione sulla
protezione dei dati che qui non è stata fatta: il dataset è pseudonimizzato ma
resta a livello di singolo evento donazionale, e la scheda \emph{Data Table} ne
consente l'esportazione integrale. \textbf{Una versione pubblica dovrebbe
esporre solo aggregati.}

% Sul piano analitico, la scheda \emph{Regression} offre tre covariate e una
% interazione, mentre la specifica del \cref{ch:glm} ne ha di più: la scheda
% serve a esplorare la direzione degli effetti, non a riprodurre il modello del
% capitolo --- il che, come mostra la \cref{sec:dash_overdisp}, non le ha
% impedito di sollevare una domanda pertinente su quel modello.

% ══════════════════════════════════════════════════════════════════════════════
%  Fine Capitolo — La Dashboard Interattiva
% ══════════════════════════════════════════════════════════════════════════════
